\documentclass[11pt,a4paper]{article}
\usepackage[utf8]{inputenc}
\usepackage{amsmath,amssymb,amsthm}
\usepackage{listings}
\usepackage{xcolor}
\usepackage[margin=1in]{geometry}
\usepackage{hyperref}

\newtheorem{theorem}{Theorem}

\lstset{
  language=C++,
  basicstyle=\ttfamily\small,
  keywordstyle=\color{blue}\bfseries,
  commentstyle=\color{green!60!black},
  stringstyle=\color{red},
  numbers=left,
  numberstyle=\tiny\color{gray},
  breaklines=true,
  frame=single,
  tabsize=4,
  showstringspaces=false
}

\title{IOI 2003 -- Amazing (Maze Exploration)}
\author{Solution}
\date{}

\begin{document}
\maketitle

\section{Problem Statement Summary}
This is an interactive problem. You are placed in an $R \times C$ grid
maze whose structure is unknown. At each step you can query which
directions (N, S, E, W) are open from the current cell. Navigate from a
start cell to a goal cell using as few moves as possible.

\section{Solution: BFS-Based Exploration}

\subsection{Strategy}
Since the maze is initially unknown, we must explore it on the fly.
We maintain a map of all explored cells and their open walls. At each
step:
\begin{enumerate}
  \item If the current cell is unvisited, query its open directions and
        record them.
  \item Run BFS on the \emph{known} graph to plan a route:
        \begin{itemize}
          \item If the goal is reachable through known cells, follow the
                shortest known path.
          \item Otherwise, navigate toward the nearest unvisited cell
                reachable from the known graph.
        \end{itemize}
  \item Execute one step of the planned route.
\end{enumerate}

\subsection{Correctness}
Every cell is visited at most once (we never re-query a visited cell).
Each navigation step either visits a new cell or moves along a shortest
path in the known graph. Since BFS always targets the nearest frontier
cell, every reachable cell is eventually explored.

\subsection{Alternative: Tremaux's Algorithm}
For a DFS-based exploration, Tremaux's algorithm guarantees visiting
each edge at most twice, giving $O(R \times C)$ total moves with
simpler bookkeeping.

\section{C++ Implementation}

\begin{lstlisting}
#include <bits/stdc++.h>
using namespace std;

const int dx[] = {-1, 1, 0, 0}; // N, S, E, W
const int dy[] = {0, 0, 1, -1};
const char dirChar[] = {'N', 'S', 'E', 'W'};

int main() {
    int R, C, sr, sc, gr, gc;
    cin >> R >> C >> sr >> sc >> gr >> gc;

    int cx = sr, cy = sc;

    auto cellKey = [](int x, int y) -> long long {
        return ((long long)x << 32) | (unsigned)y;
    };

    unordered_map<long long, vector<int>> openDirs;
    unordered_set<long long> visited;

    while (cx != gr || cy != gc) {
        long long key = cellKey(cx, cy);

        if (!visited.count(key)) {
            visited.insert(key);
            string dirs;
            cin >> dirs;
            vector<int> open;
            for (char c : dirs)
                for (int d = 0; d < 4; d++)
                    if (c == dirChar[d])
                        open.push_back(d);
            openDirs[key] = open;
        }

        // BFS to find goal or nearest unvisited cell
        queue<pair<int, int>> q;
        unordered_map<long long, long long> parent;
        unordered_map<long long, int> parentDir;
        q.push({cx, cy});
        parent[key] = -1;

        int tx = cx, ty = cy;
        bool found = false;

        while (!q.empty()) {
            auto [x, y] = q.front(); q.pop();
            long long ck = cellKey(x, y);

            if (x == gr && y == gc) {
                tx = x; ty = y; found = true; break;
            }

            if (!openDirs.count(ck)) continue;
            for (int d : openDirs[ck]) {
                int nx = x + dx[d], ny = y + dy[d];
                long long nk = cellKey(nx, ny);
                if (parent.count(nk)) continue;
                parent[nk] = ck;
                parentDir[nk] = d;
                if (!visited.count(nk) && !found) {
                    tx = nx; ty = ny; found = true;
                }
                q.push({nx, ny});
            }
        }

        // Reconstruct path from (cx,cy) to (tx,ty) and take first step
        vector<int> path;
        long long tk = cellKey(tx, ty);
        while (tk != cellKey(cx, cy)) {
            path.push_back(parentDir[tk]);
            tk = parent[tk];
        }
        reverse(path.begin(), path.end());

        if (!path.empty()) {
            cout << dirChar[path[0]] << endl;
            cx += dx[path[0]];
            cy += dy[path[0]];
        }
    }

    return 0;
}
\end{lstlisting}

\section{Complexity Analysis}
\begin{itemize}
  \item \textbf{Total moves:} $O(R \times C)$---each cell is visited at
        most a constant number of times.
  \item \textbf{Time per step:} $O(R \times C)$ for BFS over the known
        graph. Total computation: $O((R \times C)^2)$ worst case, but
        typically much better.
  \item \textbf{Space:} $O(R \times C)$ for the explored maze graph.
\end{itemize}

\end{document}
